\begin{proof}[Proof of \cref{obj:prop:fixed-alphabet-reduction}]
Fix \(d\ge2\) and write
\[
\log(ed)=1+\log d,\qquad
\alpha_d:=\frac{d}{\log(ed)}.
\]
Since \(\log d\le d-1\) and \(d>1\), we have \(0<\log(ed)\le d\), so \(\alpha_d\ge1\). Thus, for every \(n\ge1\),
\[
\frac1n
\le \min\left\{1,\frac{d}{n\log(ed)}\right\}
=\min\left\{1,\frac{\alpha_d}{n}\right\}
\le\frac{\alpha_d}{n}.
\]
% lean: CausalSmith.Stat.DiscreteBudgetvalueCurve.fixed_alphabet_reduction

The lower-bound input to \cref{obj:thm:matched-curve-frontier} is the known-distribution \(L_1\) minimax result of \citep[Theorem 3 and the unknown-both-distributions reduction following Theorem 6]{JiaoHanWeissman2018}. Applying the stated frontier bound and then the preceding comparison gives, uniformly over \(b_0\in[0,1/2]\),
\[
\frac{c_\epsilon}{n}
\le c_\epsilon\min\left\{1,\frac{d}{n\log(ed)}\right\}
\le \mathfrak R_{n,d,\epsilon,b_0}^{\mathrm{curve}}
\le C_\epsilon\min\left\{1,\frac{d}{n\log(ed)}\right\}
\le\frac{C_\epsilon\alpha_d}{n}.
\]
% lean: CausalSmith.Stat.DiscreteBudgetvalueCurve.fixed_alphabet_reduction
The risk-bound constants can be chosen as \(c_\epsilon\) and \(C_\epsilon\alpha_d\), respectively; they are positive, ordered, and fixed across \(n\) and \(b_0\).

For the identity, denote the common effect in occupied cells by \(\gamma\), so that
\[
\tau_j=\gamma>0\qquad\text{whenever }p_j>0.
\]
If \(b_0\in[0,1/2]\) and \(b\in[b_0,1-b_0]\), then \(0\le b\le1\). Cell masses satisfy \(p_j\ge0\) and \(\sum_jp_j=1\). In a zero-mass cell both sides of the following identity vanish. Hence, for every \(\pi\in\Pi_b(P)\),
\[
\sum_{j=1}^d p_j\pi_j\tau_j
=\gamma\sum_{j=1}^d p_j\pi_j
\le \gamma b,
\]
where the last inequality is the budget constraint in \cref{obj:def:budget-policy-class}.
% lean: CausalSmith.Stat.DiscreteBudgetvalueCurve.constant_effect_budget_value

Define the constant policy by \(\pi^{(b)}_j=b\) for every cell \(j\). Because \(0\le b\le1\), its coordinates are valid, and
\[
\sum_{j=1}^d p_j\pi^{(b)}_j
=b\sum_{j=1}^d p_j=b,
\qquad
\sum_{j=1}^d p_j\pi^{(b)}_j\tau_j
=\gamma b.
\]
% lean: CausalSmith.Stat.DiscreteBudgetvalueCurve.constant_effect_budget_value
Thus \(\pi^{(b)}\in\Pi_b(P)\) attains the upper bound. The maximization in \cref{obj:def:budget-value} equals \(\gamma b\), and therefore \(V_b(P)=B(P)+\gamma b\), as claimed.
\end{proof}
